\subsection{Integral parallel twists on the actual cover}
\label{sec:actual-integral-parallel-twists}

\begin{proposition}\label{prop:actual-unramified-frames}
Use the original good cover, with $A=1$, coefficient field $K$,
and level $M$.  Put
\[
       B=\{\ell:\ell\text{ prime},\ \ell\mid
                                  2d\operatorname{disc}(E)M\}.
\]
For every original component and every finite place $v$ of $K$
above $\ell\notin B$, there is an original integral framed point
on that component over a finite unramified extension of $K_v$.
Its normalized valuation index over $K_v$ is one.  In particular,
the original $\mathrm{UnramPts}(B)$ assertion holds.
\end{proposition}

\begin{proof}
The original PEL order is the maximal order $\mathcal O_E$,
its alternating lattice is self-dual, and $\ell$ divides neither
its coefficient discriminant nor its level.  The datum is of type A,
so the extra invariant $I_{\rm bad}$ of Lan's definition is one.
These facts check every good-prime condition in
\cite[Definition 1.4.1.1]{theta-twist-lan}.  The maximal-order
condition in his Condition 1.4.3.10 is automatic here.
Choose a smooth projective admissible fan and the exact original
full level completed from the half-level.  The completion is the
one already proved in the actual level construction; since $\ell$
does not divide $M$, the prescribed cyclotomic fibre is an open and
closed finite-\'etale choice.  Theorem 6.4.1.1 of the same reference
gives the original smooth proper toroidal PEL model at $\ell$,
with relative normal-crossings boundary and the tautological family.
Base change to $R=\mathcal O_{K_v}$ preserves all these properties.

In a proper smooth family, the Stein algebra is finite \'etale and
its fibres are the connected-component algebra.  The selected
geometrically integral generic component therefore extends to an
open and closed proper smooth component over $R$; its generic
constant field is already $K_v$.  It has a nonempty special fibre.
A relative normal-crossings divisor contains no component of a
fibre, so this special fibre has a closed point outside the
boundary.  Its residue field is a finite extension of the finite
residue field of $R$.  Pass to the corresponding finite unramified
extension $R'/R$.  Smoothness and Hensel's lemma lift that point to
an $R'$-point of the same abelian interior.

The lifted object retains the original polarization, coefficient
action, both labelled level maps, and primitive rigidified line.
Since $\ell\nmid\operatorname{disc}(E)$ and $K$ contains every
specified image of $E$, the coefficient algebra over $R'$ splits
as the product of the original embedding factors.  Its tangent
eigenspaces are free with precisely the prescribed ranks; choose
bases.  The primitive principal line has rank-one space of sections
over $R'$, compatible with base change, so choose a generator and
its original rigidified generic comparison.  These choices give
exactly an original integral framed point.  They are evaluations
of the finite-type framed moduli charts, hence satisfy the original
test-frame condition.  The residue field remains finite: the
source's $\mathrm{IntFr.IsAlg}$ means precisely this condition,
not algebraicity of the complete fraction field over $K$
\cite{theta-twist-source}.

The extension $R'/R$ is unramified, so restriction of its integer
valuation is the integer valuation of $K_v$.  Every evaluated rule
has value zero or an integral normalized valuation.  This is the
literal index-one criterion defining $\mathrm{UnramPts}(B)$.
\end{proof}

\begin{theorem}\label{thm:actual-short-parallel-twist}
Use the actual flat determinant sections of
Theorem~\ref{thm:actual-flat-degree}.  On each original component,
for every integer parallel shift $w=(w_\tau)_\tau$, there is a
nonzero section $t_w$ of the original parallel weight
$\operatorname{parWt}(w)$, with theta tag zero, such that
\[
 v(t_w)\ge0\quad\hbox{at every original integral test frame},
\]
and, at every complex embedding and every frame on the component,
\begin{equation}\label{eq:actual-parallel-twist-bound}
 \log\|t_w\|\le
 \frac{|p-q|}{2}|P_E|\sum_\tau|w_\tau|
       +\frac12\log\operatorname{rd}(K)
       +\sum_{\ell\in B}\log\ell+1.
\end{equation}
Both $t_w$ and its generic inverse have constant complex norm and
preserve the original cusp spaces under multiplication, with the
corresponding weight shifts.  The inverse is not asserted integral.

Under $|P_E|\le3L_b$, $\log\operatorname{rd}(E)\le3L_b$,
$\log M\le30L_b$, $\log d\le L_b$, $f=[E:\mathbf Q]/2\ge1$,
and $L_b\ge1$, the right side is at most
\[
       \tfrac32 rL_b\sum_\tau|w_\tau|
                  +\tfrac12\log\operatorname{rd}(K)+40fL_b+1.
\]
\end{theorem}

\begin{proof}
Let $s_\tau$ be the actual uncorrected determinant generator,
with the original tangent-frame norm and finite valuations.  Form
\[
                       t=\prod_\tau s_\tau^{-w_\tau}.
\]
The identity $\operatorname{parWt}(w)=
\sum_\tau(-w_\tau)\operatorname{flatWt}(\tau)$ shows that this
is the required original weight, with theta tag zero.
Let $r_\sigma>0$ be its constant complex norm and $\rho_v$ its
point-independent normalized finite valuation.  The latter have
finite support.  By Proposition~\ref{prop:actual-unramified-frames},
$\rho_v$ is an integer whenever $v$ lies above a prime outside $B$:
evaluate the same constant valuation at the index-one point.

The uncorrected degree of $s_\tau$ is $(p-q)P_E/2$; hence
\[
 D(t)=\sum_v\rho_v\frac{\log N v}{[K:\mathbf Q]}
         -\frac{1}{[K:\mathbf Q]}\sum_\sigma\log r_\sigma
       =-\frac{p-q}{2}P_E\sum_\tau w_\tau.
\]
We use the uncorrected frame here deliberately.  The correction
$q(\log d-\log\operatorname{rd}(E))$ in
$\mathrm{FlatSec}$ changes its degree bookkeeping, not the
actual evaluated section or its integral tangent lattice.

For completeness, put
$I=\prod_v\mathfrak p_v^{\lceil-\rho_v\rceil}$, a nonzero
fractional ideal.  Its logarithmic norm satisfies
\[
 \frac{\log N I}{[K:\mathbf Q]}
 \le-\sum_v\rho_v\frac{\log N v}{[K:\mathbf Q]}
                     +\sum_{\ell\in B}\log\ell.
\]
Indeed rounding costs zero outside $B$, and for each remaining
rational prime
$\sum_{v\mid\ell}\log N v\le[K:\mathbf Q]\log\ell$.
Minkowski's theorem, applied to the embedding box with radii
$R/r_\sigma$ (discs at complex places), gives $b\in I\setminus\{0\}$
with
\[
 \log(r_\sigma|\sigma b|)
 \le -D(t)+\tfrac12\log\operatorname{rd}(K)
                 +\sum_{\ell\in B}\log\ell+1
 \quad\hbox{for all }\sigma.
\]
The additive constant $1$ makes the volume strictly larger than
$2^{[K:\mathbf Q]}$ times the lattice covolume; the conjugate
radii agree because the original metric respects complex
conjugation.  This is the exact scalar construction, with the
field-degree normalization retained.

Set $t_w=bt$.  The inequalities
$v(b)\ge\lceil-\rho_v\rceil$ prove integrality at every original
test frame, and $-D(t)\le |p-q||P_E|\sum|w_\tau|/2$ proves
\eqref{eq:actual-parallel-twist-bound}.  The determinant generator
and its inverse extend as units on the generic compactification;
equivalently their original normalized norms are constant.
Multiplication by them therefore preserves holomorphy, the row
weight conditions, and cusp vanishing, in both directions.  This
argument does not require the separate theta key section.

Finally
\[
 \sum_{\ell\in B}\log\ell
 \le\log2+\log d+2f\log\operatorname{rd}(E)+\log M
 \le(6f+32)L_b\le40fL_b.
\]
\end{proof}

This supplies the actual unramified-point and integral parallel-shift
inputs.  Section~\ref{sec:actual-key-short-span} constructs the separate
key-weight shift with its finite-place normalization and applies both
shifts to the exact short-family consumer.
